\documentclass[10pt,a4paper]{amsart}
\usepackage[margin=19mm]{geometry}
\usepackage{amsmath,amssymb,mathtools}
\usepackage[scr=rsfs,cal=euler]{mathalfa}
\usepackage{xfrac}
\usepackage[unicode,hidelinks,bookmarks=false]{hyperref}
\newcommand{\ie}{i.e.\ }
\newcommand{\cf}{cf.\ }
\newcommand{\resp}{respectively\ }
\newcommand{\Subsection}{\subsection}
\providecommand{\nobreakdash}{\nobreak\hbox{-}}
\setlength{\emergencystretch}{2em}
\multlinegap0pt
\usepackage{algorithm}\usepackage{algpseudocode}
\usepackage{tikz}
\usepackage{mathtools}
\usepackage{amssymb}
\newtheorem{proposition}{Proposition}
\newcommand{\Flow}{\Phi}
\newcommand{\bfa}{\boldsymbol{a}}
\newcommand{\bfx}{\boldsymbol{x}}
\newcommand{\rd}{\mathrm d}
\title{Two-Parameter Flows for Learning Population Dynamics of Physical Systems}
\author{Paul Schwerdtner, Tobias Blickhan, Benjamin Peherstorfer}
\date{Source: arXiv:2605.26285v1 (CC BY 4.0)}
\begin{document}
\maketitle

\noindent\emph{Source and licence.} Two-Parameter Flows for Learning Population Dynamics of Physical Systems. Version arXiv:2605.26285v1 (CC BY 4.0), distributed by arXiv under the Creative Commons Attribution 4.0 International licence, http://creativecommons.org/licenses/by/4.0/, as recorded in the arXiv licence metadata for this exact version and shown on the abstract page https://arxiv.org/abs/2605.26285v1. Full licence notice: \texttt{licenses/arxiv-2605.26285-CC-BY-4.0.txt}.\par\smallskip

\noindent\textit{Editorial scope.} Counted unit: the article's proposition prop:AnalyticU (source0.tex lines 369--380). The article's complete original proof of it is delivered with it (source0.tex lines 813--851, one \texttt{proof} environment of 39 lines, which the article places away from the statement). No mathematics is written, repaired or re-derived: the statement and the proof are the article's own lines, in the article's own notation, and no retained line is substituted or re-flowed.  Retained uncounted as the source-local context the proof needs: lines 306--308; lines 350--352.  Nothing was omitted from any retained span, so the package carries no omission record.  No retained line contains a citation, so the wrapper carries no bibliography and no bibliography entry.  No retained line contains a figure environment, an image-inclusion command or drawing code, so no image is delivered and the package declares no asset dependency.  Editorial formatting only, in addition: an AMS \texttt{amsart} preamble that restates the article's own declarations and macros used by the retained lines, namely the environments \texttt{proposition} on separate counters, together with the standard packages those lines need.  The retained environments print in this document as Proposition 1 in order of appearance, whereas the article prints the same items with the numbering of its own full document.  The complete native source of the article as distributed in the arXiv e-print for this exact version is bundled unchanged as \texttt{sources/201405/source0.tex}.
\begin{equation}\label{eq:VelocityU}
        \partial_t \Flow(\bfa, t, s) = u(\Flow(\bfa, t, s), t, s)\,.
\end{equation}

\begin{align}\label{eq:SchwarzCompPDE}
        \partial_s u + v \cdot \nabla u - u \cdot \nabla v - \partial_t v = 0\,,
    \end{align}

\begin{proposition}
\label{prop:AnalyticU}
Let $\Phi$ be a $C^2$-diffeomorphism so that it is consistent and \eqref{eq:SchwarzCompPDE} holds. Set $\bfx_{t,s} = \Flow(\bfa, t, s)$ . Then, the physics-time velocity defined as
\begin{multline}
\label{eq:AnalyticUForm}
    u(\bfx_{t,s}, t, s) = D\Flow(\bfa, t, s) \\ \cdot \int_0^s \left[ D\Flow(\bfa, t, \sigma) \right ]^{-1} \partial_t v(\bfx_{t,\sigma}, t, \sigma) \, d\sigma\,.
\end{multline}
satisfies \eqref{eq:VelocityU} and its Jacobian $Du$ is given by 
\begin{equation}
    (Du)(\bfx_{t,s}, t, s) D\Flow(\bfa, t, s)  =  \frac{\mathrm d}{\mathrm d t} D\Flow(\bfa, t, s).
\end{equation}
\end{proposition}

\begin{proof}[Proof of Proposition~\ref{prop:AnalyticU}]
\emph{Analytic form of $u$:} Let $u(\bfx_{t,s}, t, s) = D\Flow(\bfa, t, s) \cdot \int_0^s \left[ D\Flow(\bfa, t, \sigma) \right ]^{-1} \partial_t v(\bfx_{t,\sigma}, t, \sigma) \, d\sigma$.

    Take derivative in s:
    \begin{align}
        \frac{\mathrm d}{\mathrm d s} u(\bfx_{t,s}, t, s) = (\partial_s u + v \cdot \nabla u)(\bfx_{t,s}, t, s)
    \end{align}
    For the right-hand side, we use:
    \begin{align}
        \frac{\mathrm d}{\mathrm d s} D\Flow(\bfa, t, s) = Dv(\bfx_{t,s}, t, s) D\Flow(\bfa, t, s)
    \end{align}
    and $\frac{\mathrm d}{\mathrm d s} \int_0^s f(\sigma) \, \mathrm d\sigma = f(s) \; \forall f$. Putting it all together, we find
    \begin{align}
        \frac{\mathrm d}{\mathrm d s} u(\bfx_{t,s}, t, s) = Dv(\bfx_{t,s}, t, s) u(\bfx_{t,s}, t, s) + \underbrace{D\Flow(\bfa, t, s) \left[ D\Flow(\bfa, t, s) \right ]^{-1}}_{= \, \mathrm{I}} \partial_t v(\bfx_{t,s}, t, s)
    \end{align}
    and hence
    \begin{align}
        \partial_s u + v \cdot \nabla u = u \cdot \nabla v + \partial_t v
    \end{align}
    evaluated at $(\bfx_{t,s}, t, s)$.
    This relation holds for all $\bfx_{t,s}$, in particular we can find for all $\bfx$ an $\bfa$ such that $\bfx = \bfx_{t,s} = \Phi(\bfa, t, s)$ at times $(t, s)$.

We have shown that $u$ satisfies the compatibility condition. It remains to show that this implies that $u \circ \Flow = \partial_t \Flow$. By using $(\partial_t \partial_s - \partial_s \partial_t) \Flow = 0$, we see that
\begin{align}
    \partial_s \left( \partial_t \Flow(\bfa, t, s) \right) = \partial_t \partial_s \Flow(\bfa, t, s) = (\partial_t v + u \cdot \nabla v)(\Flow(\bfa, t, s)) = (\partial_s u + v \cdot \nabla u)(\Flow(\bfa, t, s)),
\end{align}
using the compatibility condition in the last equality. At the same time,
\begin{align}
    \frac{\rd}{\rd s} u(\Flow(\bfa, t, s)), t, s) = (\partial_s u + v \cdot \nabla u)(\Flow(\bfa, t, s)).
\end{align}
Hence, $u \circ \Flow$ and $\partial_t \Flow$ satisfy the same linear ODE and$\frac{\rd}{\rd s} \left( u \circ \Flow - \partial_t \Flow \right) = 0$. Since $\Phi(\bfa, t, 0) = \bfa$ implies $\partial_t \Flow \big |_{s = 0} = 0$ and $u$ (defined in \eqref{eq:AnalyticUForm}) $= 0$ at $s = 0$, we conclude $u \circ \Flow = \partial_t \Flow$ on every integral curve of $\Flow$, i.e. at every $\bfx = \Flow(\bfa, t, s)$.

\emph{Jacobian of $u$:} Denote $\bfx_{t,s} = \Flow(\bfa, t, s)$. By regularity of $\Flow$,
        \begin{align}
            0 &= \left( D_{\bfa} \frac{\mathrm d}{\mathrm d t} - \frac{\mathrm d}{\mathrm d t} D_{\bfa} \right) \Flow(\bfa, t, s) 
            = D_{\bfa} u(\bfx_{t,s}, t, s) - \frac{\mathrm d}{\mathrm d t} D \Flow(\bfa, t, s)
            = (Du)(\bfx_{t,s}, t, s) D\Flow(\bfa, t, s) - \frac{\mathrm d}{\mathrm d t} D \Flow(\bfa, t, s). \nonumber
        \end{align}
\end{proof}

\end{document}
